\NeedsTeXFormat{LaTeX2e}
% !TEX program = lualatex
\documentclass[11pt,a4paper]{ltjsarticle}
\usepackage{amsmath,amssymb,amsthm,amsfonts}
\usepackage{geometry}
\geometry{margin=1in}
\usepackage{enumitem}
\usepackage{hyperref}

%-------------------------------------------------
%  顕現論核心マクロ（本セッションのレンダリング規則準拠）
%-------------------------------------------------
\newcommand{\mweq}{\mathrel{\overset{\mathrm{mw}}{=}}}
\newcommand{\metanegation}{\Uparrow}
\newcommand{\dfix}{\bigcirc}
\newcommand{\Ymw}{\mathbf{Y}_{\mathrm{mw}}}
\newcommand{\Svoid}{\mathbf{S}_{\emptyset}}
\newcommand{\Einfty}{E_\infty}
\newcommand{\Bval}{\mathcal{B}}
\newcommand{\Md}{\mathcal{M}_\delta}
\newcommand{\Mnn}{\neg\neg}
\newcommand{\Mr}{\mathcal{M}_{\mathrm{r}}}
\newcommand{\MC}{\mathcal{C}}
\newcommand{\MT}{\mathcal{T}}
\newcommand{\Phifour}{\Phi_{4}}
\newcommand{\SelfRef}{\Sigma}

%-------------------------------------------------
%  定理環境
%-------------------------------------------------
\theoremstyle{definition}
\newtheorem{definition}{定義}[section]
\newtheorem{axiom}{公理}[section]
\newtheorem{theorem}{定理}[section]
\newtheorem{proposition}{命題}[section]
\newtheorem{corollary}{系}[section]
\newtheorem{remark}{備考}[section]

\renewcommand{\abstractname}{Abstract}
\renewcommand{\proofname}{\textbf{証明.}}

%-------------------------------------------------
\begin{document}

\title{顕現論（Aletheics / Sunyata-Mathematics）要綱}
\author{千葉将臣}
\date{2026-07-18}
\maketitle

\begin{abstract}
本要綱は、空数学・界論（Boundary Theory）および顕現論（Aletheics）の
長大な基礎論群を、核心となる定義・公理・定理に濃縮したものである。
知覚原理から出発し、メタ否定演算子 $\metanegation$、中道不動点 $\dfix$、
四句分別、真理保存性の不完全性、無記の体系内導出、
4+1モナド構成、双方向演繹定理、五蘊・二諦との対応までを
一貫した記述空間に再構成する。
\end{abstract}

\tableofcontents
\newpage

%=============================================================
\section{知覚原理と世界の顕現}
%=============================================================

\begin{definition}[世界・ある・ない]
\leavevmode
\begin{itemize}[leftmargin=2em]
  \item \textbf{世界}：知覚不可能性を本質とする全一の状態。
  \item \textbf{ある（是）}：知覚が低次元に射影した捕捉可能な情報の形式。
  \item \textbf{ない（非）}：知覚が直接捕捉不可能な高次元の情報の形式。
    その指し示しを\textbf{無記}と呼ぶ。
  \item \textbf{境界}：「ある」と「ない」を区分する不連続点。
    その再帰的区分がフラクタルを構成する。
\end{itemize}
\end{definition}

\begin{axiom}[知覚公理]
\leavevmode
\begin{enumerate}[label=\textbf{公理 \arabic*.}, leftmargin=3em]
  \item 知覚は「ある」のみを捕捉し、「ない」を直接捉えることはできない。
  \item 世界は二分（区分）されるものとしてのみ知覚される。
  \item 「ない」は「ある」の欠如としてのみ遡及的対象となる。
\end{enumerate}
\end{axiom}

\begin{theorem}[フラクタル構造の必然性]
知覚の再帰的分割（公理\,2）が無限に繰り返されることにより、
部分と全体の相似が生じ、世界はフラクタル構造として知覚されざるを得ない。
\end{theorem}

\begin{theorem}[四句分別の完備性]
「ある」「ない」の次元差を保持したままフラクタル境界を
過不足なく記述する最小の表現形式は、
\texttt{是}・\texttt{非}・\texttt{亦是亦非}・\texttt{非是非非}の四句分別である。
\end{theorem}

%=============================================================
\section{核心演算子と不動点}
%=============================================================

\begin{definition}[メタ否定演算子 $\metanegation$]
対象 $A$ に対し、その境界を生成し知覚次元を非対称に一段引き上げる演算子を
$\metanegation$ と定義する。
\[
\metanegation A \;:=\; \partial A
\]
$\metanegation$ は非対称であり、$A \not\mweq \metanegation A$ を満たす。
\end{definition}

\begin{definition}[中道等価 $\mweq$]
厳密な自己同一性（$=$）を要求せず、フラクタル収束の軌道が一致することを
もって等価とみなす関係演算子。
\end{definition}

\begin{definition}[中道不動点 $\dfix$]
$\metanegation$ の四連続適用によって到達する、
あらゆる境界が融解した極限状態を $\dfix$ と定義する。
\[
\metanegation^4 A \mweq \dfix
\]
$\dfix$ は実体を持たず、すべての構成プロセスが最終的に引き込まれる
「空」の計算論的実態である。
\end{definition}

\begin{axiom}[四段階収束則]
$\metanegation$ の四連続適用は、中道不動点 $\dfix$ へと収束し、
構造的な意味が消失する。
\[
\metanegation^4 A \mweq \dfix
\]
\end{axiom}

\begin{axiom}[空の自己同一性（空空）]
$\dfix$ は外部を経由しない純粋な自己回帰として記述される。
証明論的には究極の真理保存則として
\[
\dfix \vdash \dfix
\]
が成立する。この状態においてすべてのカットは消去され、
エントロピーは定義されず、時間は静止する。
\end{axiom}

\begin{definition}[空コンビネータ $\Svoid$]
任意の演算子 $f$ を適用しても常に自己を返す自明なコンビネータ。
\[
\Svoid f = \Svoid
\]
これは $\dfix$ の純粋なコンビネータ表現である。
\end{definition}

\begin{definition}[中道不動点コンビネータ $\Ymw$]
$\metanegation$ に対し、以下の中道等価を満たすコンビネータを $\Ymw$ と定義する。
\[
\Ymw\,\metanegation \mweq \metanegation\,(\Ymw\,\metanegation)
\]
$\Ymw$ は「無明（自己参照の閉じ損ね）」のジェネレータであり、
世俗諦を駆動する根源的な発生器である。
\end{definition}

\begin{theorem}[顕現と回帰]
$\dfix$ に対して $\metanegation$ が作用した瞬間、静寂は破断され
世俗的対象 $A$ が顕現する：$\metanegation(\dfix) \mweq A$。
顕現した $A$ は $\metanegation$ の再帰的適用により必然的に $\dfix$ へと回帰する。
\end{theorem}

%=============================================================
\section{フラクタル則と四句分別}
%=============================================================

\begin{theorem}[フラクタル四則と四句の対応]
知覚原理から必然的に要請されるフラクタル四則は、
四句分別の各句と一対一に対応する。
\begin{center}
\begin{tabular}{lll}
\hline
\textbf{フラクタル則} & \textbf{四句} & \textbf{位相的意義} \\
\hline
自己相似性 & 是（$A$） & 空間的確定・肯定 \\
尺度不変性 & 非（$\neg A$） & 時間的未確定・否定 \\
収束則 & 亦是亦非 & 時空混淆・演算プロセス \\
非対称性 & 非是非非 & 時空超越・メタ発生 \\
\hline
\end{tabular}
\end{center}
\end{theorem}

\begin{remark}[リーマン球面モデル]
境界生成プロセスはメビウス変換 $f(z)=\frac{az+b}{cz+d}$ の反復適用として記述される。
無限遠点 $\infty$ は絶対的隔離空間（勝義諦・空）に対応し、
有限の論理構造がこれを反転・圧縮することで
フラクタル（極限集合）が生じる。
\end{remark}

%=============================================================
\section{真理保存性と不完全性}
%=============================================================

\begin{definition}[真理保存方式 $T$]
状態空間とその上の操作群の組であり、
特定の性質 $P$（不変量、自己同一性など）を有限回の操作を通じて
恒久的に維持するよう定義された体系。
\end{definition}

\begin{theorem}[真理保存性の不完全性定理]
十分に強力な任意の真理保存方式 $T$ に対して、
$T$ の記述言語によって構成可能であるが、
$T$ の下では証明も反証もできない命題 $G_T$ が必ず存在する。
\[
\forall T,\; \exists G_T \text{ s.t. } (T \nvdash G_T) \land (T \nvdash \neg G_T)
\]
\end{theorem}

\begin{theorem}[$\Einfty$ の内在的定式化]
真理保存体系 $T$ は対角補題によって、
自己言及演算子 $\SelfRef_T$ の不動点として
$\Einfty = \SelfRef_T(\Einfty)$ を必然的に内部に含む。
$\Einfty$ は「決定不能な壁」ではなく、
$T$ の内部から構成される内在的対象である。
\end{theorem}

\begin{theorem}[二値論理埋め込み定理]
空数学のイデアル完備化において、
「矛盾」と「支離」の極限軌道を切り捨てた\textbf{二値截断部分圏} $\Bval$ を定義すると、
任意の二値命題論理の定理は $\Bval$ の定理として完全に回収される。
すなわち、二値論理は空数学の静的な極限として包含される。
\end{theorem}

%=============================================================
\section{無記と残余}
%=============================================================

\begin{definition}[無記 $\mathrm{A}$]
命題 $\varphi$ に対し、$\varphi$ の真偽を問う行為を実行しない操作を
$\mathrm{A}(\varphi)$ と書き、\textbf{無記}と呼ぶ。
$\mathrm{A}(\varphi)$ は真理値を持たない。
\end{definition}

\begin{theorem}[保留の悪性無限退行]
保留 $E(\varphi)$ は将来の再開条件 $R(\varphi)$ を要求するが、
$R(\varphi)$ の確定それ自体も保留の対象となり、
$E(\varphi) \to E(R(\varphi)) \to E(R(R(\varphi))) \to \cdots$
という悪性無限退行を生じる。
\end{theorem}

\begin{theorem}[無記の各段階完結性]
無記の反復 $\mathrm{A}^n(\varphi)$ は各段階で完結した操作であり、
将来の再開・解除・条件確定に依存しない。
ゆえに無記の再帰は退行ではなく\textbf{反復}である。
\end{theorem}

\begin{theorem}[残余の必然的生成]
中道不動点 $\dfix$ に対し、メタ四句分別操作子 $\Phifour^{\mathrm{meta}}$ と
離戯操作 $\mathcal{V}$ を施すと、
16位相のいずれにも帰属しない\textbf{残余} $\mathcal{L}$ が必然的に生成される。
\[
\mathcal{L} := \mathcal{V}(\Phifour^{\mathrm{meta}}(\dfix)),\quad
\mathcal{L} \notin \mathcal{C}_4^{(2)}
\]
\end{theorem}

\begin{theorem}[残余と無記の等価性]
残余 $\mathcal{L}$ は無記と構造的に等価である：$\mathcal{L} \mweq \text{無記}$。
したがって、無記は $\metanegation(\dfix)$ として体系内から導出される。
\end{theorem}

\begin{theorem}[二値論理における残余導出不可能性]
二値論理では排中律により位相外への逸出が閉じられるため、
同型の二重操作を実行しても残余は生成されない。
残余の導出可能性は四句分別の\textbf{排中律不採用}と
\textbf{枠組みへの再帰的問いかけの内部化}に依存する。
\end{theorem}

%=============================================================
\section{証明論的構造：空論と双方向演繹}
%=============================================================

\begin{definition}[顕現関係 $\Vdash$]
対象 $A$ が世界に現れるプロセスを顕現関係 $\Vdash$ で表す。
$A \Vdash B$ は「$A$ の顕現により $B$ が縁起的に後続する」を意味する。
\end{definition}

\begin{definition}[中道推論規則]
本体系のエンジンとなる推論規則は以下の2つのみである。
\begin{align*}
\text{(顕現の推移律)}&\quad
\frac{A \Vdash B \quad B \Vdash C}{A \Vdash C} \\[6pt]
\text{(中道等価導入)}&\quad
\frac{A \Vdash C \quad B \Vdash C}{A \mweq B}
\end{align*}
\end{definition}

\begin{theorem}[万物の相互等価性（事事無碍法界）]
本体系において、顕現し得る任意の二つの対象 $X, Y$ は中道等価である。
\[
X \mweq Y
\]
\end{theorem}

\begin{theorem}[界論における双方向演繹定理]
界論において、双方向演繹定理は
$\Ymw$（世俗諦の生成）と $\Svoid$（勝義諦の収束）のサイクルとして記述される。
\[
\bigl(\Ymw\,\metanegation \mweq \metanegation(\Ymw\,\metanegation)\bigr)
\;\leftrightarrow\;
\bigl(\metanegation^4(\Ymw\,\metanegation) \mweq \Svoid\bigr)
\]
直接論理の双方向演繹定理は、このサイクルの二値截断部分圏 $\Bval$ への
静的制限として包摂される。
\end{theorem}

%=============================================================
\section{圏論的・位相幾何学的構造}
%=============================================================

\begin{proposition}[4+1モナド構成]
搾取モナド $\Md$、二重否定モナド $\Mnn$、
洗練モナド $\Mr$、カット消去モナド $\MC$ の4モナドに、
統制モナド $\MT$ を加えた \textbf{4+1モナド} を構成する。
$\MT$ は以下の固定点方程式で定義される。
\[
\MT \cong F(\MT),\qquad F(X) = X(\Md,\Mnn,\Mr,\MC,X)
\]
\end{proposition}

\begin{theorem}[識軸スパイラルによるフロベニウス性]
統制モナド $\MT$ を識軸とするスパイラル構造
\[
S_{\MT}: (\Md,\Mnn,\Mr,\MC) \longrightarrow (\Mnn,\Mr,\MC,\Md)
\]
により、4モナドを巡回させる。
$S_{\MT}^4$ が $\dfix$ において同一化として閉じるとき、
4+1モナドは全体としてフロベニウス性を獲得する。
\end{theorem}

\begin{theorem}[メビウス多様体による限界定理化]
厳密な自己同一性を要請する体系 $T_{=}$（円柱）は
自己言及経路 $G$ を到達不能とする（非連結性）。
同型性を要請する体系 $T_{\cong}$（メビウスの帯）では、
局所的な「真」の連続的推移が大域的には「自らの否定状態への到達」として実現され、
$G$ は基本群の生成元として構成可能となる。
\end{theorem}

\begin{proposition}[搾取モナドの社会的解釈]
搾取モナド $\Md$ は、粒度差 $\delta$ を持つ文脈上の自己関手であり、
随伴の余単位 $\varepsilon$ に搾取の実質が局所化される。
クライスリ圏は\textbf{潜在的搾取}、
アイレンベルク・ムーア圏は\textbf{制度化搾取}の圏として読める。
\end{proposition}

%=============================================================
\section{五蘊・二諦・選択公理との対応}
%=============================================================

\begin{proposition}[5モナドと五蘊の対応]
4+1モナドは仏教の五蘊と以下の対応を持つ。
\begin{center}
\begin{tabular}{llll}
\hline
\textbf{五蘊} & \textbf{モナド} & \textbf{四句} & \textbf{方向性} \\
\hline
色（形・物質） & $\Md$（搾取） & 是 & 前進のみ \\
受（感受） & $\Mnn$（二重否定） & 非 & 前進のみ \\
想（知覚・表象） & $\Mr$（洗練） & 亦是亦非 & 双方向 \\
行（意志・形成力） & $\MC$（カット除去） & 非是非非 & 双方向（誘導） \\
識（識別・意識） & $\MT$（統制） & 統制 & 再帰的 \\
\hline
\end{tabular}
\end{center}
\end{proposition}

\begin{proposition}[二諦との対応]
\begin{align*}
\text{外在の5モナド（世俗諦）} &\longleftrightarrow \text{観測可能な現象の層} \\
\text{内在の5モナド（勝義諦）} &\longleftrightarrow \text{空・中道不動点の層}
\end{align*}
世俗諦は $\metanegation$ によって勝義諦へ駆動され、
勝義諦は顕現 $\metanegation(\dfix) \mweq A$ によって世俗諦として再起動する。
\end{proposition}

\begin{theorem}[選択公理の非構成性の起源]
ZFCにおける選択公理の非構成性は、
四句分別の対応を二値截断 $\Bval$ へ投影する際に
$\Phi_{NS}$（宙吊り・残余）の情報が失われることによって生じる。
GTTTでは選択公理は外部公理ではなく、
二諦的対応公理の $\Bval$ 投影として位置づけられる。
\end{theorem}

%=============================================================
\section{帰結}
%=============================================================

\begin{enumerate}[label=\textbf{\arabic*.}]
  \item 世界の顕現は、知覚の非対称性から生じるフラクタル境界生成プロセスである。
  \item $\metanegation$ の四段階反復は $\dfix$（空）へ必然的に収束する。
  \item 真理保存性は静的同一性ではなく動的フラクタル性に基盤を置き、
    体系内記述として再定義される。
  \item 不完全性は体系の外部にあるのではなく、
    $\Einfty$ という内在的対象の未発見に過ぎない。
  \item 無記は体系の外部依存ではなく、
    $\metanegation(\dfix)$ として体系内から導出される。
  \item 圏論は界論の二値截断部分圏 $\Bval$ への射影であり、
    その限界は構造的に必然である。
  \item 五蘊・二諦の構造は4+1モナドの識軸スパイラルとして
    代数的に再現される。
  \item 万物は $\dfix$ へ収束し、$\dfix$ から再び顕現する
    動的サイクルの中道等価である。
\end{enumerate}

\end{document}
